\documentclass[11pt]{article}
\usepackage[margin=1in]{geometry}
\usepackage{amsmath,amssymb}
\usepackage[hidelinks]{hyperref}

\newcommand{\D}{\mathbb D}
\newcommand{\T}{\mathbb T}

\begin{document}

\section*{Human Review: Pringsheim-singular spaceability question}

\noindent Original automatic proof: \texttt{solution\_packet.pdf}

\noindent Checked date: 13 June 2026

\noindent Status: Manually checked.

\noindent Verdict: Not a solution to the original question as stated in
Bernal-Gonzalez--Bonilla--Lopez-Salazar--Seoane-Sepulveda.

\section*{Human Comments}

The proposed construction appears to prove a weaker natural-boundary
statement, but it does not prove the Pringsheim-singularity condition required
in the original question.

The source packet interprets ``Pringsheim singular at every boundary point''
as meaning that the holomorphic function has no analytic continuation through
any point of \(\T\). Under this interpretation, Hadamard's gap theorem is a
natural tool: the constructed functions have lacunary Taylor support and
radius of convergence exactly \(1\), so the unit circle is a natural boundary.

However, the original paper
\[
\text{Bernal-Gonzalez--Bonilla--Lopez-Salazar--Seoane-Sepulveda,}
\]
\emph{Nowhere holderian functions and Pringsheim singular functions in the
disc algebra}, defines Pringsheim singularity for the boundary function in a
stronger real-variable sense. If
\[
  f_0(t)=f(e^{it}),
\]
then \(f|_{\T}\) is Pringsheim-singular at \(e^{it_0}\) when the Taylor series
of \(f_0\) at \(t_0\) has radius of convergence zero. Equivalently,
\[
  \limsup_{r\to\infty}
  \left|
    \frac{f_0^{(r)}(t_0)}{r!}
  \right|^{1/r}
  =
  +\infty.
\]
This is stronger than merely saying that \(f\) has no holomorphic continuation
through \(e^{it_0}\).

\section*{Where the Proof Falls Short}

The packet proves that every nonzero element of the constructed closed span
has \(\T\) as a natural boundary. This rules out analytic continuation through
any boundary point, and plausibly gives a nowhere-real-analytic boundary
function. But it does not establish the derivative-growth condition above.

For a boundary function
\[
  f_0(t)=\sum_{n\in\Lambda} c_n e^{int},
\]
one has
\[
  f_0^{(r)}(t_0)
  =
  \sum_{n\in\Lambda} (in)^r c_n e^{int_0}.
\]
To prove Pringsheim singularity in the sense of the original paper, one would
need pointwise lower bounds on these derivative sums at every \(t_0\), strong
enough to force
\[
  \left|f_0^{(r)}(t_0)/r!\right|^{1/r}
\]
to be unbounded.

The coefficient choice \(c_n=e^{-\sqrt n}\) is suggestive: a single term of
the form \(n^r e^{-\sqrt n}\) can be very large compared with \(r!\) for
appropriate \(n\). But the proof does not show that such a term dominates the
whole derivative sum, nor does it rule out cancellations. This is especially
important for general nonzero elements of the closed span, where several
lacunary blocks can contribute simultaneously with complex coefficients.

\section*{Review Outcome}

This packet should be rejected as a solution to the original Pringsheim
spaceability question. It should not be archived as a correct full answer to
the question in \([6]\).

It may still be useful as a proof of a weaker statement: spaceability of
smooth disk-algebra functions whose unit circle is a natural boundary, or
possibly whose boundary restrictions are nowhere real analytic. But the
argument does not prove that the boundary Taylor radius is zero at every
point, which is the condition required for Pringsheim singularity in the
original reference.

\end{document}
